% Desarrollo basado en los apuntes del estudiante; no es una transcripción.
\par\textbf{\textcolor{unabblue}{\faBrain}\quad La paradoja de la IA generativa: potencia técnica y vacío emocional}\par
La proliferación de modelos generativos en la producción audiovisual y la comunicación digital reduce drásticamente los costes marginales de síntesis y los tiempos de desarrollo. Sin embargo, esta aceleración revela una paradoja fundamental: la optimización del proceso técnico no garantiza resonancia humana ni eficacia comunicativa. Cuando la tecnología se asume como un fin en sí misma y no como un medio instrumental, las narrativas devienen estériles. Para la Ingeniería en Sistemas, esto evidencia que la capacidad de cómputo carece de impacto duradero si no responde a las motivaciones, vulnerabilidades y necesidades reales de las personas.

\subsubsection*{\textcolor{unabblue}{\faHeart}\quad La «verdad del consumidor» como ancla del diseño y la comunicación}
En la convergencia entre tecnología, mercadotecnia y comportamiento humano, el concepto de \textit{verdad del consumidor} (\textit{consumer insight}) constituye el núcleo de toda narrativa efectiva: una revelación psicológica o tensión cultural no evidente que explica cómo los individuos perciben su entorno e interactúan con él.

\begin{itemize}
  \item \textbf{Límites algorítmicos:} Los modelos probabilísticos detectan correlaciones en grandes volúmenes de datos históricos, pero son incapaces de vivenciar empatía, arraigo o afecto.
  \item \textbf{Imperativo de conexión:} El receptor final es un ser humano guiado por juicios emocionales e intuitivos; si el mensaje carece de autenticidad humana, se cataloga de inmediato como ruido artificial.
\end{itemize}

En el diseño de sistemas y productos interactivos, la tecnología debe amplificar la verdad del usuario, no enmascarar la falta de empatía o propósito en la propuesta de valor.

\clearpage
\subsubsection*{\textcolor{unabblue}{\faFilm}\quad Estudio de caso: la nostalgia de Coca-Cola (1995 vs. recreación con IA)}
Un caso emblemático para examinar los límites de la generación sintética frente a la conexión emocional lo ofrece la comunicación navideña de la compañía Coca-Cola:

\begingroup
\small
\renewcommand{\arraystretch}{1.25}
\begin{tabularx}{\linewidth}{@{}>{\raggedright\arraybackslash}p{2.8cm}>{\raggedright\arraybackslash}X>{\raggedright\arraybackslash}X@{}}
\toprule
\textbf{Dimensión} & \textbf{Campaña tradicional (1995)} & \textbf{Recreación sintética con IA} \\
\midrule
\textbf{Producción y técnica} & Grabación física con actores, caravana real de camiones iluminados y dirección cinematográfica en set. & Generación algorítmica completa mediante modelos de difusión de video, sin rodaje físico. \\
\textbf{Percepción de la audiencia} & Evocación de calidez, magia comunitaria, asombro infantil y un ritual festivo compartido. & Sensación de artificio, distanciamiento, expresiones inorgánicas y efecto de «valle inquietante». \\
\textbf{Impacto en la marca} & Consolidación de un símbolo cultural perdurable y refuerzo del vínculo afectivo. & Crítica masiva en redes y medios por percibir un ahorro operativo a costa del alma y la autenticidad. \\
\bottomrule
\end{tabularx}
\endgroup

El anuncio de 1995 (\textit{Holidays Are Coming}) basaba su éxito en la imperfección orgánica de la realidad: miradas genuinas de asombro en los niños, el reflejo tangible de luces de tungsteno sobre la nieve y la calidez del contacto interpersonal. En contraste, la versión generada íntegramente mediante algoritmos generativos detonó una reacción de rechazo casi unánime en el público.

La audiencia no cuestionó la resolución cromática ni la capacidad computacional requerida para sintetizar las tomas; rechazó la evidente ausencia de humanidad. Las figuras humanas generadas por IA evidenciaban movimientos no naturales y miradas desprovistas de emoción genuina, provocando la sensación de un producto manufacturado en serie para abaratar costes en lugar de celebrar una tradición comunitaria. Este fenómeno demuestra que el usuario no busca la perfección matemática en una historia, sino la verdad afectiva que la sustenta.

\clearpage
\subsubsection*{\textcolor{unabblue}{\faCompass}\quad El debate ético y creativo: ¿qué merece ser contado y por qué?}
El dilema central que plantea la inteligencia artificial contemporánea trasciende la eficiencia operativa y exige formular una pregunta deontológica: \textit{¿qué merece ser contado y por qué?}

Herramientas avanzadas de lenguaje y generación multimedia (como Claude, plataformas generativas de difusión o modelos multimodales) constituyen instrumentos excepcionales de ejecución, prototipado y experimentación. Reducen fricciones operativas y democratizan el acceso a la producción visual compleja. No obstante, la delegación del criterio editorial, la intencionalidad ética y la sensibilidad social en un algoritmo representa una abdicación del rol crítico del profesional.

\begin{enumerate}
  \item \textbf{La técnica responde al «cómo»:} Los modelos computacionales resuelven con rapidez la sintaxis visual, la textura gráfica y la coherencia formal de un producto.
  \item \textbf{La persona responde al «qué» y al «por qué»:} Solo el ser humano posee la capacidad de discernir el significado, evaluar las consecuencias éticas de un mensaje y garantizar que la narrativa construida aporte valor real a la sociedad.
\end{enumerate}

Para los ingenieros en sistemas y creadores tecnológicos, la automatización jamás puede constituir una justificación para diluir el mensaje humano. La verdadera innovación no reside en reemplazar a las personas en los espacios donde se construyen el significado y la empatía, sino en emplear la tecnología para potenciar nuestra capacidad de conectar, comunicar y transformar positivamente el entorno social.
